\originalchapter{氢原子的量子化}
\originalsection{两体约化与尺度}
氢原子发光谱有离散频率, 提示内部能量差离散. 现在将中心势方法用于解释束缚能量与位置分布; 质心的自由运动能量不计为内部能级. 电子和原子核都可运动. 设质量为$m_e,M$, 质心坐标为$\boldsymbol X=(m_e\boldsymbol r_e+M\boldsymbol r_N)/(m_e+M)$, 相对坐标为$\boldsymbol r=\boldsymbol r_e-\boldsymbol r_N$. 对只依赖相对坐标的相互作用, 动能分为质心项和相对项:
\[
H=\frac{P^2}{2(m_e+M)}+\frac{p^2}{2\mu}-\frac{g}{r},
\quad \mu=\frac{m_eM}{m_e+M},\quad g=\frac{Ze^2}{4\pi\varepsilon_0}.
\]
$Z$为核电荷整数,$e>0$为基本电荷大小, 核$+Ze$、电子$-e$,$\varepsilon_0$为SI真空介电常数. $P=p_e+p_N$共轭于质心$X$, 相对共轭动量$p=(Mp_e-m_ep_N)/(m_e+M)$. 反解$p_e=m_eP/(m_e+M)+p,p_N=MP/(m_e+M)-p$, 平方相加交叉项抵消, 就得到上式动能. $r$是两粒子间距, 不假定核在实验室固定. 这里使用SI库仑常数; 官方作业的Gaussian单位中写$g=Ze^2$. 两种写法不能混用. 对氢$Z=1$, 且$M\gg m_e$, 所以$\mu$接近$m_e$, 但谱线精密计算不能永远忽略差别.

定义长度$a_Z=\hbar^2/(\mu g)$, 它使动能尺度$\hbar^2/(\mu a_Z^2)$和势能尺度$g/a_Z$相当. 这个量纲估计说明原子大小的来源, 还没有得到精确能级. 精确量子化须解径向方程和两端边界.

\originalsection{径向级数与能量}
$E=0$是无限远分离阈值, 束缚态$E<0$. 远处$u\prime\prime\simeq\kappa^2u$, 所以$\kappa^{-1}$为指数尾衰减长度; 用$\rho=2\kappa r$把这个尺度移除, $\eta$则衡量库仑项强度. 束缚态$E<0$, 设$\kappa=\sqrt{-2\mu E}/\hbar$, $\rho=2\kappa r$, $\eta=\mu g/(\hbar^2\kappa)$. 式\eqref{eq:radial}变为
\[
\frac{\mathrm d^2u}{\mathrm d\rho^2}
+\left[-\frac14+\frac\eta\rho-\frac{\ell(\ell+1)}{\rho^2}\right]u=0.
\]
在零附近, 标准三维Coulomb实现取$u(0)=0$, 代$u\sim\rho^s$到最奇异两项给$s(s-1)=\ell(\ell+1)$, 两根$\ell+1,-\ell$, 正则解为$\rho^{\ell+1}$; $\ell=0$时另一支常数虽然半轴平方可积, 仍被此边界排除. 在无穷处, 可积解含$\e^{-\rho/2}$. 因而写$u=\rho^{\ell+1}\e^{-\rho/2}F(\rho)$. 令$g_0=\rho^{\ell+1}\e^{-\rho/2}$, 则$g_0\prime/g_0=(\ell+1)/\rho-1/2$, $u\prime\prime=g_0[F\prime\prime+2(g_0\prime/g_0)F\prime+(g_0\prime\prime/g_0)F]$, 代入并除掉公共因子后,
\[
\rho F''+(2\ell+2-\rho)F'+(\eta-\ell-1)F=0.
\]
设$F=\sum_{s\ge0}c_s\rho^s$, 比较$\rho^s$系数得到
$c_{s+1}=(s+\ell+1-\eta)c_s/[(s+1)(s+2\ell+2)]$.

先看递推的实际用法: $c_1=(\ell+1-\eta)c_0/(2\ell+2)$, $c_2=(\ell+2-\eta)c_1/[2(2\ell+3)]$. 例如$\ell=0,\eta=1$给$F=1$, $\ell=0,\eta=2$给$F\propto1-\rho/2$, $\ell=1,\eta=2$又给$F=1$. 特殊函数增长定理在这里作为明确的引用工具: 该方程原点正则解若不在有限阶终止, 无穷端含$\e^\rho$增长支, 因而乘$\e^{-\rho/2}$后仍不正规. 它通常名为合流超几何函数的大$\rho$渐近性质, 本册不让读者先学复Gamma函数再理解量子化. 避免这一项需要分子在某个$s=n_r$处为零, 即$\eta=n_r+\ell+1=:n$. 因而$n=1,2,\ldots$, $\ell=0,1,\ldots,n-1$, 且
\[
E_n=-\frac{\mu g^2}{2\hbar^2n^2},\qquad
u_{n\ell}\propto r^{\ell+1}\e^{-r/(na_Z)}
L_{n-\ell-1}^{2\ell+1}\!\left(\frac{2r}{na_Z}\right).
\]
$n_r=n-\ell-1$是径向多项式次数, 也为内部径向节点数, $n$标内部能量,$\ell$标角分布,$m$标轴投影. 这里$L$表示关联Laguerre多项式; 多项式形式正是有限递推的结果. 非终止解的增长判据可由上述方程的合流超几何大参数渐近验证, 不是仅凭“大$s$像指数”对任意级数作一般定理. 正则起点和可积终点共同排除其他能量.

\originalsection{基态分布和维里关系}
氢样原子基态为$\psi_{100}=(\pi a_Z^3)^{-1/2}\e^{-r/a_Z}$. 位置密度$|\psi|^2$在$r=0$最大, 但半径落在$[r,r+\dd r]$的概率为$4r^2\e^{-2r/a_Z}\dd r/a_Z^3$. 对其取导, 得最可能的\emph{半径}是$a_Z$. 密度最大位置与半径分布最大值不同, 原因是球壳体积随$r^2$增长.

用$\int_0^\infty r^s\e^{-\lambda r}\dd r=s!/\lambda^{s+1}$, 得$\langle r\rangle=3a_Z/2$, $\langle r^2\rangle=3a_Z^2$, $\langle1/r\rangle=1/a_Z$. 若核半径$R\ll a_Z$, 核内概率约为$4R^3/(3a_Z^3)$, 忽略指数的相对误差为$O(R/a_Z)$.

这也解释光谱: 忽略精细结构,$n=2\to1$能差$13.6(1-1/4)=10.2$ eV, 发出光子$\hbar\omega=10.2$ eV; 还需选择定则判断具体态是否允许. 伸缩生成元使位置与动量反向缩放, 表示压窄原子会提高动能而增强吸引势的平衡. 取伸缩生成元$D=(\boldsymbol r\cdot\boldsymbol p+\boldsymbol p\cdot\boldsymbol r)/2$. 在适当束缚态定义域,$[D,H]=\ii\hbar(2T-\boldsymbol r\cdot\nabla V)$. 定态中任意无显含时算符的期望值不变, 所以$2\langle T\rangle=\langle\boldsymbol r\cdot\nabla V\rangle$. 对$V=-g/r$, 有$\boldsymbol r\cdot\nabla V=g/r=-V$, 因而$\langle T\rangle=-\langle V\rangle/2=-E$. 此关系检验符号和能量尺度, 不证明电子沿一条Bohr圆轨道运动.
